\chapter{イントロダクション}

力学では、質点の運動$\gamma:I\to\mathbb{R}^3$の軌跡を求めることが最も重要な目標でした。
そのために運動方程式$F(t,\gamma(t))=m\gamma''(t)$から様々な$\gamma$を求めていきました。

ところで、この曲線の性質があらかじめ最初から分かっていて、それらを用いて運動方程式はごく最小限に用いて運動の様子が分かればうれしいところです。
それを抜きにしても、運動する質点の描く軌跡は神秘に満ちていて、調査しがいのあるというものです。

今回はそんな動機や魅力にとらわれた一人である、ガウスの曲線論を展開していきましょう。



\section{微分幾何学とは何か}

ガウスの曲線論は、現代\textbf{微分幾何学}と呼ばれる分野の走りとなります。
微分可能な写像$\gamma:I\to\mathbb{R}^3$に対して、微分を駆使して曲がり具合（曲率）や、ねじれ具合（捩率）といった量を定義して、不変量などを調べます。
微分幾何は曲線論に限らず、\textbf{曲面論}、あるいはさらなる高次元版の\textbf{リーマン幾何}など多岐にわたっています。
いずれにしても、何か計算しやすい量を少ない道具から導き、その性質をよく調べるということをします。

なぜこれらを調べるかというと、これらの「計算しやすい対象」の性質を知っておくことで、曲線がどのような形をしているかある程度推察できるからです。
この講義では、最終的には曲率と捩率から曲線が定まってしまうという「曲線論の基本定理」を一つの目標とします。
余裕があれば、これらを使って物理学や工学などへの応用について述べます。
例えば高速道路や新幹線のレールを建てるときに、変な方向にGがかかる設計だと車体が高速で吹き飛んでしまいます。
そういったことが起こらないように、曲線論で見積もることができるのです。
そんなことをやっていきましょう。
